\chapter{Сигнал тревоги}
% полная версия: chapters/14_pain.tex

\pencilfig{14}{\pencilart{14}}{Боль --- пожарная сигнализация, а не термометр.}

Все остальные чувства сообщают, каков мир. Боль --- нет. Она сообщает, что что-то угрожает телу, и её громкость во многом задаёт сама машина, а не раздражитель. Боль --- это пожарная сигнализация, а не термометр.

Датчики боли срабатывают только на сильное: на жар выше сорока с лишним градусов, на сильное давление, на едкие вещества. Они привыкают гораздо слабее, чем датчики прикосновения, --- сигнализация не должна замолкать, пока угроза есть.

\section*{Два провода}

Вспомните, как вы ударились пальцем. Сначала --- острая, точная боль: где и когда. Через мгновение --- тупая, разлитая, долгая: насколько плохо. Это два разных провода. Один быстрый, его сигнал долетает за десятки миллисекунд. Другой медленный, его сигнал идёт порядка секунды.

\section*{Ворота}

Почему, ударившись, мы трём ушиб? В спинном мозге сигнал боли проходит через «ворота». Прикосновение будит тормозной нейрон, который прикрывает их. Поэтому потереть ушиб действительно помогает --- но только против умеренной боли; сильная проходит всё равно. Что сильная боль при этом сама выключает тормозной нейрон, --- часть исходной теории ворот, её прямого подтверждения пока нет.

\pencilfig{126}{%
% gate: the pain wire drives the relay neuron; touch wakes an inhibitory neuron that closes the gate
\draw[pen=pcAmber] (0,1.35) -- (3.05,1.0); \draw[arr=pcAmber] (2.8,1.03) -- (3.08,1.0);
\node[lbl, text=pcAmber!70!black, anchor=east] at (-0.05,1.35) {боль};
\draw[pen=pcBlue] (0,-0.15) -- (1.75,-0.15); \draw[arr=pcBlue] (1.5,-0.15) -- (1.78,-0.15);
\node[lbl, text=pcBlue, anchor=east] at (-0.05,-0.15) {прикосновение};
\draw[dotpen=pcRed, wash=pcRed] (2.05,-0.15) circle (0.26);
\draw[pcRed, line width=0.7pt, -{Bar[width=6pt]}] (2.25,0.05) -- (3.2,0.66);
\node[lbl, text=pcRed, anchor=north] at (2.05,-0.45) {тормоз};
\draw[dotpen=pcGray, wash=pcGray] (3.4,0.9) circle (0.34);
\node[lbl, anchor=south] at (3.4,1.3) {ворота};
\draw[pen] (3.75,0.9) -- (5.6,0.9); \draw[arr] (5.35,0.9) -- (5.65,0.9);
\node[lbl, anchor=west] at (5.7,0.9) {в мозг};
}{Ворота в спинном мозге. Сигнал боли проходит в мозг через нейрон-ворота, а прикосновение будит тормозной нейрон, который их прикрывает. Поэтому потереть ушиб помогает, но только против умеренной боли.}

\section*{Ручка громкости сверху}

Мозг умеет сам делать боль громче или тише. Сверху, из головного мозга, к воротам идут радиоканалы --- \nt{СЕРО}, \nt{НОРА} и особые вещества, похожие на обезболивающие, которые мозг выделяет сам. В опасной ситуации боль может почти исчезнуть: бегущему от опасности некогда хромать. Ожидание облегчения тоже приглушает боль, отчасти через те же вещества.

\section*{Когда сигнализация залипает}

Повторная боль повышает чувствительность: тот же удар начинает болеть сильнее. Это тоже обучение --- петля, которая сама себя подкручивает, пока рана заживает. Но у сильной петли есть опасность: она может «залипнуть» и продолжать звонить, когда рана уже зажила, --- как группа нейронов из второй главы, которая поддерживает сама себя.

\section*{Учитель и прерывание}

Боль --- ещё и учитель: для дофаминового «лучше или хуже, чем ожидалось» она работает как отрицательная награда. И она прерывает текущее дело, как вспышка норадреналина: бросай всё, займись опасностью. А самая быстрая реакция --- отдёрнуть руку от горячего --- замыкается прямо в спинном мозге той самой парой мышц-противников из прошлой главы, ещё до того, как вы успели почувствовать боль.

\fullversion{14}
